% Solutions of Book 4, chapter 23 (context-free grammars and pushdown automata).
\section*{Chapter \ref{ch:b2:context-free-languages} --- Context-Free Grammars and Pushdown Automata}

\begin{solution}{exo:b2:context-free-languages:1}
\omoutput{code/bachelor-2/23-context-free-languages/outputs/exo_leftmost.txt}{The \omterm{def:b2:context-free-languages:parse-tree}{leftmost derivation} of $(a) * a$.}
The tree: $E$ has one child $T$; that $T$ has the children $T$, $*$, $F$; the left $T$ goes down through $F$ to $(E)$,
whose $E$ goes down through $T$ and $F$ to $a$; the right $F$ is $a$. Its nine inner nodes are the eight steps plus the
\omterm{def:b2:context-free-languages:grammar}{start symbol} $E$.
\end{solution}

\begin{solution}{exo:b2:context-free-languages:2}
\omoutput{code/bachelor-2/23-context-free-languages/outputs/exo_brackets.txt}{The words of length at most 4, and a
check up to length 10.}
$S \Rightarrow (S)S \Rightarrow ()S \Rightarrow ())S(S \Rightarrow ())(S \Rightarrow ())($. The grammar generates the
words with as many $($ as $)$. Each production adds one of each, so every \omterm{def:b2:context-free-languages:derivation}{generated} word has equal numbers. Conversely,
let $w$ be nonempty with equal numbers, say starting with $($, and let $(u)$ be its shortest nonempty prefix with equal
numbers: the difference between the numbers of $($ and $)$ stays positive inside it, so it ends with $)$, and $u$ and
the rest $v$ both have equal numbers; by induction on the length they are \omterm{def:b2:context-free-languages:derivation}{generated}, and $S \to (S)S$ gives $w$. A
word starting with $)$ is handled by $S \to )S(S$ in the same way.
\end{solution}

\begin{solution}{exo:b2:context-free-languages:3}
\omoutput{code/bachelor-2/23-context-free-languages/outputs/exo_palindromes.txt}{The grammar against every word.}
$P \to aPa \mid bPb \mid a \mid b \mid \varepsilon$: a palindrome of length at least 2 starts and ends with the same
letter around a shorter palindrome, and those of length 0 and 1 are the last three \omterm{def:b2:context-free-languages:grammar}{productions}. There are $2^{\lceil
n/2 \rceil}$ palindromes of length $n$, and $1 + 2 + 2 + 4 + 4 + \cdots + 64 + 64 = 253$ up to length 12.
\end{solution}

\begin{solution}{exo:b2:context-free-languages:4}
Five trees, one per way to bracket the three additions: $((a + a) + a) + a$, $(a + (a + a)) + a$, $(a + a) + (a + a)$,
$a + ((a + a) + a)$ and $a + (a + (a + a))$; this is $C_3 = 5$. The grammar with levels gives one, the first, since $E
\to E + T$ puts the last addition at the root (\cref{out:b2:context-free-languages:ambiguity}).
\end{solution}

\begin{solution}{exo:b2:context-free-languages:5}
\omoutput{code/bachelor-2/23-context-free-languages/outputs/exo_cyk.txt}{The two tables.}
For $aabb$: the letters give $A, A, B, B$; the middle $ab$ is $S$ by $S \to AB$; $abb$ is $C$ by $C \to SB$; the whole
word is $S$ by $S \to AC$. For $abab$ the two $ab$ are $S$, but nothing derives $aba$, $bab$ or the whole word. $S \to
AB$ gives $ab$ and $S \to AC \to aSb$ wraps a word of the language in $a \cdots b$: the language is $\{a^nb^n : n \ge
1\}$.
\end{solution}

\begin{solution}{exo:b2:context-free-languages:6}
\omoutput{code/bachelor-2/23-context-free-languages/outputs/exo_unit_cycle.txt}{The naive recursion, and the method.}
Following $A \to B$ then $B \to A$ leads back to $A$, and the recursion never ends (here it is stopped at depth 12;
otherwise the stack overflows). The method searches the graph of the unit \omterm{def:b2:context-free-languages:grammar}{productions} and visits each \omterm{def:b2:context-free-languages:grammar}{nonterminal} at
most once, so from $S$ it reaches $\{S, A, B\}$ and gives $S$ the \omterm{def:b2:context-free-languages:grammar}{productions} $ab$, $a$ and $b$; $A$ and $B$ are then no
longer reachable from $S$ and are dropped.
\end{solution}

\begin{solution}{exo:b2:context-free-languages:7}
\omcode{code/bachelor-2/23-context-free-languages/exercises/descent.rs}{9}{36}{One function per \omterm{def:b2:context-free-languages:grammar}{nonterminal}}
\omoutput{code/bachelor-2/23-context-free-languages/outputs/exo_descent.txt}{The recogniser against the grammar.}
The two \omterm{def:b2:context-free-languages:grammar}{productions} of $S$ start with different letters, $c$ or the $a$ of $T$, so one letter decides; $T$ has one
production. Each function returns whether it read one word of its \omterm{def:b2:context-free-languages:grammar}{nonterminal}, moving the position past it, and the
whole word is accepted when $S$ succeeds and nothing is left. Every word of the language has a length $3k + 1$.
\end{solution}

\begin{solution}{exo:b2:context-free-languages:8}
Every \omterm{def:b2:context-free-languages:derivation}{generated} word is balanced, by induction on the \omterm{def:b2:context-free-languages:derivation}{derivation}: $\varepsilon$ is, and if $u$ and $v$ are, so is
$(u)v$: its numbers are equal, a prefix inside $(u$ has one more $($ than a prefix of $u$, and a prefix $(u)v'$ has as
many of each as $v'$. Conversely, by strong induction on the length: let $w$ be balanced and nonempty. Its first letter
is $($, since the prefix of length 1 cannot have more $)$. Let $(u)$ be the shortest nonempty prefix with as many $)$
as $($; inside it the difference stays positive, so $u$ is balanced, and the rest $v$ is balanced too (its prefixes
add to a balanced prefix). Both are shorter than $w$, so $S \Rightarrow^* u$ and $S \Rightarrow^* v$, and $S
\Rightarrow (S)S \Rightarrow^* (u)v = w$.
\end{solution}

\begin{solution}{exo:b2:context-free-languages:9}
\omcode{code/bachelor-2/23-context-free-languages/listings/exo_pda_equal.rs}{9}{16}{Six moves on the surplus, and an
accepting move}
\omoutput{code/bachelor-2/23-context-free-languages/outputs/exo_pda_equal.txt}{The automaton on every word up to length
10.}
The stack holds the surplus: $k$ symbols $A$ above the bottom if $k$ more $a$ than $b$ have been read, $k$ symbols $B$
if $k$ more $b$. An $a$ pushes an $A$ unless a $B$ is on top, which it pops; symmetrically for $b$. By induction on
the input read, the stack is always that surplus, so the bottom is on top exactly when the numbers are equal, and the
last move accepts then. This automaton guesses only when the counts are equal and input remains: it may read on or
take the spontaneous move to the \omterm{def:b2:finite-automata:dfa}{accepting state}, and the wrong choice dies there, which the search over runs
handles. A deterministic version keeps ``counts equal'' in its state: two states, equal (accepting) and unequal,
moving from one to the other when the bottom symbol comes back on top or leaves it.
\end{solution}

\begin{solution}{exo:b2:context-free-languages:10}
Let $L = \{ww\}$ and $R = a^+b^+a^+b^+$. If $L$ were \omterm{def:b2:context-free-languages:derivation}{context-free}, so would be $L \cap R = \{a^ib^ja^ib^j : i, j \ge
1\}$ (\cref{exo:b2:context-free-languages:12}). Take $p$ from \cref{thm:b2:context-free-languages:pumping} and $z =
a^pb^pa^pb^p$. Since $|vwx| \le p$, $vwx$ lies inside one block or across two adjacent ones. If $v$ or $x$ contains two
different letters, $uv^2wx^2y$ is not in $R$. Otherwise pumping with $i = 2$ makes one or two adjacent blocks longer
and leaves the others unchanged: the first and third blocks, or the second and fourth, no longer have the same length.
Either way $uv^2wx^2y \notin L \cap R$, a contradiction.
\end{solution}

\begin{solution}{exo:b2:context-free-languages:11}
\omcode{code/bachelor-2/23-context-free-languages/exercises/bigcount.rs}{80}{108}{CYK with counts}
\omoutput{code/bachelor-2/23-context-free-languages/outputs/exo_bigcount.txt}{Sums of 10 to 100 operands.}
A number is a vector of 32-bit digits; a product of two digits plus a digit and a carry stays below $2^{64}$, so each
step is exact in 64 bits. Each cell maps a \omterm{def:b2:context-free-languages:grammar}{nonterminal} to its number of trees: one per production $A \to a$ at the
bottom, and for each split and each $A \to BC$ the product of the counts of $B$ and $C$, added to $A$'s count. In
\omterm{def:b2:context-free-languages:cnf}{Chomsky normal form} a tree is exactly a choice of production and split at each node, so the sum counts each tree once;
here the conversion of $E \to E + E \mid a$ (into $E \to E\,X_0$, $X_0 \to \text{\code{'+'}}\,E$) keeps the trees in
bijection. The Catalan numbers are computed by $C_k = C_{k-1} \cdot 2(2k - 1)/(k + 1)$, a division that is always
exact. $C_{99}$ has 57 digits.
\end{solution}

\begin{solution}{exo:b2:context-free-languages:12}
Let $P$ be a PDA with states $Q$ and \omterm{def:b2:finite-automata:dfa}{accepting states} $F$, and $D$ a complete DFA with states $R$, start $r_0$,
transitions $\delta$ and \omterm{def:b2:finite-automata:dfa}{accepting states} $F'$. The product PDA has states $Q \times R$, start $(q_0, r_0)$, the same
stack, \omterm{def:b2:finite-automata:dfa}{accepting states} $F \times F'$, and the moves $((q, r), a, Z) \to ((q', \delta(r, a)), \gamma)$ for each move
$(q, a, Z) \to (q', \gamma)$ of $P$ reading a letter, and $((q, r), \varepsilon, Z) \to ((q', r), \gamma)$ for each
spontaneous move. By induction on the number of moves, the product goes from $((q_0, r_0), w, Z_0)$ to $((q, r), y,
\gamma)$ exactly when $P$ goes from $(q_0, w, Z_0)$ to $(q, y, \gamma)$ and $r = \delta^*(r_0, x)$ with $w = xy$. So
it accepts $w$ exactly when $P$ accepts it and $D$ does. A product of two PDAs would need two stacks, which is why
the intersection of two \omterm{def:b2:context-free-languages:derivation}{context-free languages} can fail to be \omterm{def:b2:context-free-languages:derivation}{context-free}.
\end{solution}

\begin{solution}{pb:b2:context-free-languages:1}
\textbf{1.} $D \Rightarrow D\,D \Rightarrow [D]\,D \Rightarrow [w]\,D \Rightarrow [w]w$.

\textbf{2.} The nonempty sequences of items, an item being a word $w$, a nonempty notice in bold between two $*$, or
a nonempty notice as a link between brackets.

\textbf{3.} Intersected with the \omterm{def:b2:regular-languages:regular}{regular language} $[^*w\,]^*$, the notices give $\{[^nw\,]^n\}$; for $i \ne j$ the
words $[^i$ and $[^j$ have different residuals ($w\,]^i$ is in one only), so this is not regular, and since \omterm{def:b2:regular-languages:regular}{regular
languages} are closed under intersection, neither are the notices.

\textbf{4.} $D \to wD \mid w$.

\textbf{5.} Both have the root $D \to D\,D$: one with $ww$ under the left child, $(ww)w$, the other with $ww$ under the
right child, $w(ww)$.

\textbf{6.} One bold item \code{*(w *w* w)*}, whose inner sequence of three items is grouped in two ways; and the
sequence \code{*w*}, \code{w}, \code{*w*}, grouped in two ways.

\textbf{7.} $D \to D\,D$, which is associative but says nothing of the grouping. With $D \to ID \mid I$, a sequence is a
first item followed by the rest, in one way only.

\textbf{8.} The list of three items \code{*w*}, \code{w}, \code{*w*}; and one item \code{*D*} whose $D$ is the list \code{w},
\code{*w*}, \code{w}.

\textbf{9.} With the same symbol for opening and closing, a $*$ can be read either way, which question 8 shows. With
\code{<} and \code{>} a reader always knows which it is, as with brackets.

\textbf{10.} That no notice of length at most 8 has two trees, nothing more. A proof for every length: in $I$, the
first letter chooses the production; in $D \to ID \mid I$, the item ends where its closing symbol matches its opening
one (or after one $w$), so the split is forced, and an induction on the length gives one tree.

\textbf{11.} Yes: both generate the nonempty sequences of items with bold between $*$; only the trees differ.

\textbf{12.} \code{D} $\Rightarrow$ \code{I D} $\Rightarrow$ \code{<D> D} $\Rightarrow$ \code{<I> D} $\Rightarrow$ \code{<w> D} $\Rightarrow$ \code{<w> I} $\Rightarrow$ \code{<w>[D]} $\Rightarrow$ \code{<w>[I]} $\Rightarrow$ \code{<w>[w]}.

\textbf{13.} TERM (the four symbols in the right sides of length 3), BIN (cutting \code{<D>} and \code{[D]} into
\code{'<'}~$X_0$, $X_0 \to D$~\code{'>'} and the same with $X_1$) and UNIT ($D \to I$). No \omterm{def:b2:context-free-languages:grammar}{nonterminal} derives
$\varepsilon$, so START and DEL change nothing.

\textbf{14.} Because of the unit production $D \to I$: every item is also a notice.

\textbf{15.} Each letter's cell holds the \omterm{def:b2:context-free-languages:grammar}{nonterminals} with a production to it: \code{'<'} for \code{<}, $D$ and $I$ for
$w$ (by $I \to w$, and $D \to w$ from UNIT), \code{'['}, \code{']'} and \code{'>'} for the three others.

\textbf{16.} It covers $[w]$: $X_1$ derives $w]$ by $X_1 \to D$~\code{']'} (the cell of length 2 at $w$), and
\code{'['}~$X_1$ is a right side of both $I$ and $D$.

\textbf{17.} The top cell, of length 7, is empty, so $D$ is not in it. The factor \code{[w>} is the trouble: no production
puts a \code{>} right after a link's content, so no cell over a factor containing it holds $D$, $I$, $X_0$ or $X_1$ ending
there correctly.

\textbf{18.} $(7^3 - 7)/6 = 56$ and $(6^3 - 6)/6 = 35$.

\textbf{19.} As in \cref{thm:b2:context-free-languages:cyk}: $\sum_{\ell=2}^{n} (n - \ell + 1)(\ell - 1) = \sum_{j=1}^{n-1}
(n - j)j = (n^3 - n)/6$.

\textbf{20.} The two columns agree on all five lines; for $n = 21$, $(9\,261 - 21)/6 = 1\,540$.

\textbf{21.} About 7 to 8 times: the work is cubic, and $2^3 = 8$; for the smallest notices the time is too short to
measure well.

\textbf{22.} Seven: $D \to ID$, $D \to$ \code{'<'}~$X_0$, $D \to$ \code{'['}~$X_1$, $I \to$ \code{'<'}~$X_0$, $I \to$
\code{'['}~$X_1$, $X_0 \to D$~\code{'>'}, $X_1 \to D$~\code{']'}; each split point tests all seven, so the time is
about $7 \times (n^3 - n)/6$ tests.

\textbf{23.} Scaling the last line by $(20\,000/336)^3 \approx 211\,000$: about $0.078 \times 211\,000 \approx
16\,000$ seconds, some four and a half hours. Checking a long notice needs a linear-time method for a restricted kind
of grammar, which is what the parsers of compilers use.

\textbf{24.} \emph{CYK's cubic count}: on a grammar in \omterm{def:b2:context-free-languages:cnf}{Chomsky normal form}, CYK examines $(n^3 - n)/6$ split points
for a word of length $n$, each against every production $A \to BC$, so its time grows as $n^3$; the program counted
exactly that number for notices of 21 to 336 letters, and its time grew about eightfold at each doubling.
\end{solution}
